\section{Computations}\label{app:comp}

This appendix specifies the finite computations of the paper. Every result marked as computer-assisted in the text is the output of a terminating exact search over a finite set. Section~\ref{app:method} describes the enumeration of commuting block maps, Section~\ref{app:proofs} the computations on which computer-assisted propositions rely, Section~\ref{app:stated} the computations whose results are stated in remarks, and Section~\ref{app:confirm} computations that confirm results proved by hand. The programs accompany the paper as supplementary material, in its directory \texttt{code/}; the file \texttt{README.md} there states, for each program, what it computes and its expected output.

\subsection{Enumeration of commuting block maps}\label{app:method}

\emph{Commutation as a finite constraint.} Let $h$ be a block map of radius $R$ with local rule $h_{\mathrm{loc}}$, and $F=F_f$ with $f$ elementary. Both $h\circ F$ and $F\circ h$ are block maps of radius $R+1$, so they are equal if and only if their local rules agree on every word $w\in\A^{2R+3}$. Write $n=2R+1$ and, for $w=(w_0,\dots,w_{2R+2})$, let $v_1,v_2,v_3$ be the three subwords of length $n$ starting at positions $0,1,2$, and let $v_F$ be the word of length $n$ obtained by applying the local rule of $F$ to $w$. Then
\[
(h\circ F)(w)=h_{\mathrm{loc}}(v_F),\qquad (F\circ h)(w)=\Bigl(h_{\mathrm{loc}}(v_2)_y,\ h_{\mathrm{loc}}(v_2)_x\oplus f\bigl(h_{\mathrm{loc}}(v_1)_y,h_{\mathrm{loc}}(v_2)_y,h_{\mathrm{loc}}(v_3)_y\bigr)\Bigr),
\]
where the subscripts $x,y$ denote the two output bits. Writing $X_v,Y_v$ for the two unknown bits of the table entry $h_{\mathrm{loc}}(v)$, each word $w$ contributes the two constraints
\begin{equation}\label{app:eq:constraints}
X_{v_F}=Y_{v_2},\qquad Y_{v_F}\oplus X_{v_2}=f(Y_{v_1},Y_{v_2},Y_{v_3}).
\end{equation}
The first is two binary clauses; the second is a constraint on at most five variables, encoded by its truth table in at most sixteen clauses. Altogether there are $2\cdot4^{2R+1}$ variables and at most $18\cdot4^{2R+3}$ clauses. For an automaton of radius $r$ of the form \eqref{gen:eq:F} the same encoding is used with words of length $2R+2r+1$ and the $2r+1$ subwords of length $n$ that $f$ reads.

\emph{Enumeration and verification.} All solutions are enumerated by repeated solving with blocking clauses, using PySAT \cite{PySAT} with CaDiCaL \cite{CaDiCaL} or MiniSat \cite{MiniSat}. Every table returned by the solver is re-checked, independently of the clauses, by evaluating $h\circ F$ and $F\circ h$ on all words. Each solution is then matched against the maps $\sh^aF^b$ of radius at most $R$, which are built directly by composing local rules, and against the constant maps. The enumeration asserts that every $\sh^aF^b$ of radius at most $R$ occurs among the solutions, which checks the encoding against the maps known in advance to commute.

\emph{Invertibility.} A block map is an automorphism of $\A^{\Z}$ if and only if it is injective. Injectivity of a block map of radius $R$ is decided exactly by the pair graph: its vertices are the pairs $(u,v)$ of words of length $2R$, with an edge $(u,v)\to(u',v')$ labelled $(a,b)\in\A^2$ whenever $u'$ and $v'$ are $ua$ and $vb$ with their first letters removed, and $h_{\mathrm{loc}}(ua)=h_{\mathrm{loc}}(vb)$. Then $h$ fails to be injective if and only if a bi-infinite path passes through an off-diagonal vertex, which is decided by iteratively deleting vertices without a predecessor or without a successor. Where only the invertible commuting maps are wanted, balance is added to the clauses: a surjective block map of radius $R$ is balanced \cite{Hedlund1969}, each output symbol occurring exactly $4^{2R}$ times in its table, so the balanced commuting maps are enumerated and each is submitted to the pair-graph test. The supplementary material contains \texttt{sat/centralizer.py} for the elementary rules, \texttt{sat/cent\_r.py} for the automata \eqref{gen:eq:F} of radius $r$, and \texttt{sat/invertible\_sat.py} for the enumeration with balance.

\subsection{Computations used in computer-assisted propositions}\label{app:proofs}

\emph{The certificate $(\mathrm T_P)$.}\label{app:general} For given $r$, $g$ and $P$, the sets $\mathrm{Im}_m$ of Section~\ref{gen:sec:tower} are computed by iterating the successor relation on the $2^{2rP}$ column states until they stabilise at $S_\infty$, and $(\mathrm T_P)$ is then tested. Increasing $P$ gives a semi-decision procedure for a sufficient condition; we do not claim that \textup{($\Lambda$)} implies $(\mathrm T_P)$ for some $P$.

\emph{Proposition~\ref{gen:prop:r2TP}.} For each of the $65\,536$ functions $g$ of four variables, $(\mathrm T_P)$ is tested for $P=1,2,3,4$ in turn, up to the least $P$ for which it holds. The supplementary material contains two independent programs for this, which represent the column states differently (\texttt{scan\_r2.py} with \texttt{scan\_r2b.py}, and \texttt{tp\_direct.py}), and a third (\texttt{compare.py}) that compares their tables of least $P$ for every $g$; the tables are identical. The program \texttt{scan\_r2.py} also builds $\Lambda_2$ and $\Lambda_3$ from the cell rule for the functions with no $(\mathrm T_P)$, $P\le3$: for the $44$ functions with least $P=4$ the graphs $\Lambda_2$ and $\Lambda_3$ are strongly connected, and for $g=\xi_3$ and $g=1\oplus\xi_3$ the graph $\Lambda_3$ is not. For $g=\xi_3$ the graphs $\Lambda_M$, $2\le M\le5$, are built in the same way and are not strongly connected. The program \texttt{xcheck.py} compares the two $(\mathrm T_P)$ routines on $153$ functions and counts the functions that satisfy \textup{(N)} and \textup{(D)}.

\emph{Tables~\ref{gen:tab:r2}, \ref{gen:tab:r2ten} and~\ref{gen:tab:r3}, and Propositions~\ref{gen:prop:r2} and~\ref{gen:prop:r3}.} $J(g)$ and self-duality are computed from the truth table. $(\mathrm T_P)$ is tested for $P\le4$ at radius $2$ and $P\le3$ at radius $3$. Condition $(\mathrm I)_q$ is tested by enumerating, for each $n$, the points of least spatial period $n$, computing the stabiliser of each $\Z^2$-orbit, and counting the orbits with each stabiliser; this gives $k(\mathcal L)$ exactly (Lemma~\ref{gen:lem:stab}). The witness listed is the first in the order of increasing $n$, then increasing index $nb$, then increasing $s$; the enumeration goes up to $n=5$. Each of these quantities is computed by two independent programs, with identical results; the supplementary material contains \texttt{examples.py} and \texttt{checks3.py}. In addition, $\Lambda_M$ built from the cell rule is strongly connected for $M\le3$ for every rule of Tables~\ref{gen:tab:r2} and~\ref{gen:tab:r2ten}, and for $M\le2$ for every rule of Table~\ref{gen:tab:r3}. The column ``strip'' of Table~\ref{gen:tab:r2} is a SAT encoding of \eqref{gen:eq:star}, as a system of constraints on the truth table of $\gamma$, whose enumeration stops after $20$ solutions (\texttt{sat/strip.py}); the column ``extra'' is the enumeration of Section~\ref{app:method} (\texttt{sat/cent\_r.py}). The statements about rules $90$ and $150$ in Section~\ref{gen:sec:examples} are computed in the same way.

\subsection{Computer-assisted statements in remarks}\label{app:stated}

\emph{Affine rules} (Remark~\ref{pm:rem:linear}). $\Lambda_M$ and the eventual image of $\Gamma_{w,0}$ are built directly from the local rule and tested for strong connectivity. For rule $90$ the graph $\Lambda_M$ is not strongly connected for $1\le M\le8$; this is the range quoted in Remark~\ref{pm:rem:linear} and in Section~\ref{gen:sec:examples}.

\subsection{Confirmations of results proved by hand}\label{app:confirm}

The following computations confirm results that the paper proves by hand. None of them is used in a proof.

\emph{Rule $30$ at radius at most $4$} (Theorem~\ref{thm:rigidity30}). The enumeration of Section~\ref{app:method} for $f=30$ and $R=1,2,3,4$ gives exactly $2R^2+2R+3$ tables, that is $7$, $15$, $27$ and $43$: the $2R^2+2R+1$ tables of the maps $\sh^aF_{30}^b$ with $|a|+|b|\le R$, which are pairwise distinct and invertible, and the two constant maps onto $(0,0)^{\Z}$ and $(1,1)^{\Z}$. So $\Cent_4(F_{30})$ consists of the elements of $\langle\sh,F_{30}\rangle$ of radius at most $4$ and these two constants, in agreement with Theorem~\ref{thm:rigidity30}; this also determines the non-invertible commuting block maps of radius at most $4$, which the theorem does not address. The enumeration is part of \cite{PaperD}.

\emph{The elementary permutive rules} (Theorem~\ref{pm:thm:main}). For $f=45,75,135$ and $R=1$, and for $f=45,135$ and $R=2$, the enumeration gives only the maps $\sh^aF^b$ and constant maps. For the eight non-affine left-permutive rules, $\Lambda_M$ built from the cell rule is strongly connected for $M\le5$.

\emph{Rules that read two neighbours} (Section~\ref{sec:halfradius}). The three finite checks, Lemmas~\ref{hr:lem:P}, \ref{hr:lem:sections} and~\ref{hr:lem:level1}, are confirmed by direct computation (\texttt{halfradius/hr\_checks.py}): the section identities by free reduction of the section words and by evaluating both sides on all words of length $7$. The window graph $W_M$, built from the cell rule, is strongly connected for $M\le6$, and the level-$1$ image of $\{t|_{x_*^M}:t\in G,\ t(x_*^M)=x_*^M\}$ is exactly $A_4$ for $119$ when $1\le M\le10$ (and $S_4$ for $M=0$), and $S_4$ for $187$ when $0\le M\le10$. An exhaustive SAT enumeration (\texttt{halfradius/kernel\_sat.py}) of the right-sided block maps of anticipation at most $6$ that commute with $F$ and are bijective in their first argument, as every right-sided automorphism with right-sided inverse is (Lemma~\ref{hr:lem:onesided}(3)), gives exactly the powers $F^b$ with $|b|\le6$, for both $119$ and $187$. For rule $102$ the same enumeration gives non-invertible right-sided block maps commuting with $F$. The enumeration of \cite{PaperD} gives, for the classes of $63$ and $187$, exactly the $(2R+1)^2$ automorphisms $\sh^aF^b$ of radius at most $R$ for $R\le3$ (Remark~\ref{hr:rem:census}).
